% Part: lambda-calculus
% Chapter: lambda-definability
% Section: lambda-definable-recursive

\documentclass[../../../include/open-logic-section]{subfiles}

\begin{document}

\olfileid[tr]{lam}{dfl}{ldr}
\olsection{\usetoken{S}{lambda definable} Fonksiyonlar Özyinelemelidir}

Bütün kısmi özyinelemeli fonksiyonlar !!{lambda definable} olmakla kalmaz;
tersi de doğrudur. Yani bütün !!{lambda definable} fonksiyonlar kısmi
özyinelemelidir.

\begin{thm}
  \ollabel{thm:lambda-computable} Bir kısmi $f$ fonksiyonu
  !!{lambda definable} ise kısmi özyinelemelidir.
\end{thm}

\begin{proof}
  İspatın yalnızca taslağını veriyoruz. Önce $\lambd$-terimlerini
  aritmetikleştiririz; yani dizilerin alışılmış asal kuvvetleri kodlamasını
  kullanarak $\lambd$-terimlerine dizgesel biçimde Gödel sayıları atarız.
  Sonra bir lambda teriminin Gödel sayısı $t$ üzerinde çalışan kısmi
  özyinelemeli $\fn{normalize}(t)$ fonksiyonunu tanımlarız: terimin normal
  biçimi varsa bunun Gödel sayısını döndürür, yoksa tanımsızdır. Ayrıca doğal
  sayılar ile bunlara karşılık gelen Church sayılarının Gödel sayıları
  arasında iki yönde eşleme yapan kısmi özyinelemeli $\fn{toChurch}$ ve
  $\fn{fromChurch}$ fonksiyonlarını tanımlarız.

  Bu özyinelemeli fonksiyonları kullanarak $f$'yi kısmi özyinelemeli bir
  fonksiyon olarak tanımlayabiliriz. $f$'yi !!{lambda define}s bir
  $\lambd$-terimi $F$ vardır. $f(n_1,\dots,n_k)$'yi hesaplamak için önce
  $\fn{toChurch}(n_i)$ ile karşılık gelen Church sayılarının Gödel sayılarını
  elde ederiz; bunları $\Gn{F}$'ye ekleyerek
  $F\num{n_1}\dots\num{n_k}$ teriminin Gödel sayısını kurarız. Şimdi bu
  Gödel sayısına $\fn{normalize}$ uygularız. $f(n_1,\dots,n_k)$ tanımlıysa
  $F\num{n_1}\dots\num{n_k}$ bir normal biçime sahiptir ve bu normal biçim
  bir Church sayısı olmak zorundadır. $f(n_1,\dots,n_k)$ tanımsızsa normal
  biçim yoktur; dolayısıyla
  \[
  \fn{normalize}(\Gn{F\num{n_1}\dots\num{n_k}})
  \]
  tanımsızdır. Son olarak normalleştirilmiş terimin Gödel sayısına
  $\fn{fromChurch}$ uygularız.
\end{proof}

\end{document}
